\section{Πειραματική Διάταξη και Υπολογιστικό Περιβάλλον}
\label{sec:eval_setup}

\subsection{Υποδομή Κατανεμημένου Συστήματος}
Τα πειράματα εκτελέστηκαν σε πλήρως απομονωμένο και αναπαραγώγιμο περιβάλλον κατανεμημένων εικονικών εμπορευματοκιβωτίων (Docker containers), διασυνδεδεμένων μέσω του εικονικού δικτύου γέφυρας \texttt{simulated-lan}. Η κατανεμημένη τοπολογία περιλαμβάνει:
\begin{itemize}
    \item \textbf{Κόμβοι LEAD DHT ($N=3$):} Τρεις ανεξάρτητοι φυσικοί/λογικοί κόμβοι υλοποιημένοι σε Rust, οι οποίοι διαχειρίζονται από κοινού 60 εικονικούς κόμβους (vnodes, 20 vnodes ανά κόμβο) επί δακτυλίου τύπου Chord.
    \item \textbf{Νησιωτικοί Οδηγοί Orchestrator ($M=2\dots 4$):} Αυτόνομοι εξελικτικοί κινητήρες γραμμένοι σε Rust (χρησιμοποιώντας τον ασύγχρονο εκτελεστή \texttt{Tokio} και το πλαίσιο \texttt{Tonic} gRPC), που διαχειρίζονται ανεξάρτητους υπο-πληθυσμούς και επικοινωνούν μέσω δακτυλίου μετανάστευσης.
    \item \textbf{Μικροϋπηρεσία Υποκατάστασης (Surrogate Node):} Κόμβος ταχείας πρόβλεψης (Rust με C++/HIP engine σε AMD ROCm GPU), ο οποίος παρέχει surrogate inference και επιγραμμική επανεκπαίδευση (online retraining) στο παρασκήνιο.
    \item \textbf{Συστοιχία Υπολογιστικών Εργατών (Worker Pool):} Κλιμακούμενος αριθμός εργατών Python/AeroSandbox ($W \in \{1, 2, 4, 8, 16, 32\}$), καθένας εκ των οποίων εκτελεί πλήρη αεροδυναμική προσομοίωση 3D Vortex Lattice Method (VLM).
    \item \textbf{Διαμεσολαβητής Εξισορρόπησης Φορτίου Envoy:} Reverse proxy υψηλής επίδοσης που δρομολογεί τα αιτήματα αξιολόγησης Tier 3 προς τους εργάτες με πολιτικές Round-Robin και Least-Request (P2C).
    \item \textbf{Δίαυλος Τηλεμετρίας και Παρακολούθησης:} Συλλογή δομημένων αρχείων καταγραφής JSON μέσω Fluent-Bit, ροή συμβάντων μέσω Apache Kafka (KRaft mode) και απεικόνιση σε πραγματικό χρόνο μέσω Node.js WebSocket dashboard.
\end{itemize}

\subsection{Γεωμετρικό Μοντέλο Αεροσκάφους και 10-Διάστατο Γονιδίωμα}
Το προς βελτιστοποίηση αεροσκάφος αποτελεί ένα τρισδιάστατα εκτυπώσιμο μη επανδρωμένο εναέριο όχημα (UAV). Ο σχεδιαστικός χώρος αναπαρίσταται από ένα συνεχές 10-διάστατο διάνυσμα γονιδίων $\mathbf{x} \in [0, 1]^{10}$, το οποίο κανονικοποιείται γραμμικά στα φυσικά όρια σχεδιασμού:
\begin{enumerate}
    \item \textbf{Εκπέτασμα πτέρυγας ($b$):} $b \in [200, 600]\text{ mm}$ (ημι-εκπέτασμα $b/2$).
    \item \textbf{Χορδή ρίζας ($c_{\text{root}}$):} $c_{\text{root}} \in [40, 100]\text{ mm}$.
    \item \textbf{Χορδή ακροπτερυγίου ($c_{\text{tip}}$):} $c_{\text{tip}} \in [20, 60]\text{ mm}$ ($c_{\text{tip}} \le c_{\text{root}}$).
    \item \textbf{Γωνία βέλους ($\Lambda$):} $\Lambda \in [0^\circ, 25^\circ]$.
    \item \textbf{Γωνία δίεδρου ($\Gamma$):} $\Gamma \in [0^\circ, 10^\circ]$.
    \item \textbf{Γωνία συστροφής ($\theta_{\text{twist}}$):} $\theta_{\text{twist}} \in [-5^\circ, 2^\circ]$.
    \item \textbf{Διαμήκης θέση πτέρυγας ($X_{\text{wing}}$):} Θέση έδρασης κατά μήκος της ατράκτου.
    \item \textbf{Αεροτομή πτέρυγας (NACA 4-digit):} Παραμετρική αεροτομή τύπου NACA MPTT, όπου $M$ είναι η μέγιστη καμπυλότητα ($0\dots 6\%$), $P$ η θέση της ($20\dots 50\%$ της χορδής) και $T$ το σχετικό πάχος ($8\dots 16\%$).
    \item \textbf{Ουραίο πτέρωμα (Οριζόντιο \& Κατακόρυφο):} Αυτόματος δυναμικός υπολογισμός διαστάσεων μέσω όγκων ουράς (Stan Hall tail volume sizing) με συμμετρική αεροτομή NACA 0010 για εγγύηση στατικής διαμήκους ευστάθειας ($C_{m_\alpha} < 0$).
    \item \textbf{Εσωτερικός ωφέλιμος όγκος ($V_{\text{fuse}}$):} Αυστηρός γεωμετρικός περιορισμός $V_{\text{fuse}} \ge 20.000\text{ mm}^3$ για τη στέγαση μπαταρίας και ηλεκτρονικών πτήσης (avionics).
\end{enumerate}

Στο Σχήμα~\ref{fig:wing_geometry_3d} παρουσιάζεται η τρισδιάστατη γεωμετρική απόδοση της πτέρυγας, η κατανομή πάνελ VLM και η αεροδυναμική στρωτή ροή του βελτιστοποιημένου αεροσκάφους. Στον Πίνακα~\ref{tab:ga_optimization_results} παρατίθενται οι φυσικές διαστάσεις και τα αεροδυναμικά χαρακτηριστικά σύγκρισης μεταξύ της αρχικής βασικής διαμόρφωσης (baseline) και της εξελιγμένης λύσης.

\begin{figure}[htbp]
    \centering
    \includegraphics[width=0.95\textwidth]{figures/fig1_wing_geometry_3d.png}
    \caption{Τρισδιάστατη γεωμετρική επικύρωση αεροσκάφους: (α) 3D απεικόνιση πτέρυγας, ατράκτου και ουραίου πτερώματος με πάνελ VLM, (β) κατανομή αεροτομών NACA κατά το εκπέτασμα και (γ) ισοζύγιο αεροδυναμικών δυνάμεων σε κατάσταση ισόρροπης πτήσης cruise.}
    \label{fig:wing_geometry_3d}
\end{figure}

../../evaluation/figures/table1_geometric_aero_specs.tex

Όπως φαίνεται στον Πίνακα~\ref{tab:ga_optimization_results}, η εξελικτική διαδικασία πέτυχε αύξηση της αεροδυναμικής απόδοσης $L/D$ κατά $+17.8\%$ (από $11.07$ σε $13.04$), ενώ παράλληλα μείωσε το μήκος διαδρόμου απογείωσης κατά $53.4\%$ χάρη στη βελτιωμένη άντωση κατά τη χαμηλή ταχύτητα.
